\chapter{학습 계획표}

\begin{summarybox}{한눈에}
\textbf{기관 운영의 핵심 메뉴}. 학생별 또는 반별로 ``무슨 콘텐츠를 언제까지 풀어야 하는지'' 매트릭스로 관리하고, 학생 제출물·글쓰기·테스트 결과까지 한 화면에서 추적합니다. 5개 탭으로 구성되어 있어요.
\end{summarybox}

\kfShot{08-study-plans}{학습 계획표 대시보드 — 5개 탭}

\section*{화면 구조}

\screentabs{\activetab{캘린더} \quad \activetab{학생별} \quad \inacttab{제출물} \quad \inacttab{글쓰기} \quad \inacttab{테스트}}

\begin{screenbox}{학생별 매트릭스 (가장 많이 쓰는 화면)}
\textbf{유료 학생의 학습 계획표}

\medskip
\noindent
\begin{tabular}{@{}p{1.5cm}p{2cm}p{2cm}p{2cm}p{2cm}p{2cm}@{}}
\toprule
\textbf{범위} & \textbf{국어농장 1} & \textbf{국어농장 2} & \textbf{학습활동} & \textbf{테스트} & \textbf{글쓰기} \\
\midrule
기본 & 미수행 & 배정 전 & 미제출 & — & 미제출 \\
2단원 & 완료 & 미수행 & — & 응시 전 & — \\
3 & 배정 전 & 배정 전 & — & — & 미제출 \\
\bottomrule
\end{tabular}

\medskip
\textbf{좌측 상단}: 기관 / 반 / 학생 드롭다운 \quad
\textbf{우측 상단}: 🖨 일괄 PDF 인쇄
\end{screenbox}

\section*{여기서 할 수 있는 일}
\begin{itemize}[leftmargin=1.2em]
  \item \textbf{캘린더 탭} — 일자별 배정·마감 일람
  \item \textbf{학생별 탭} — 매트릭스 (행=범위/단원, 열=콘텐츠 타입). 셀 클릭 시 모달 → 학습 결과 / 콘텐츠 배정 / PDF 인쇄 버튼
  \item \textbf{제출물 탭} — 학생이 제출한 학습활동 (사진·PDF) 검수
  \item \textbf{글쓰기 탭} — 학생 작문 제출물 (지식과 지혜 와 별도로 학습 계획표에 묶인 글쓰기)
  \item \textbf{테스트 탭} — 배정한 시험의 응시·채점 현황
  \item \textbf{콘텐츠 배정} — 셀 클릭 → ``콘텐츠 배정'' → 검색 모달 (4개 탭: 종합 검색 / 농장별 모드 / 프로 모드 / 내용 숙지) → 콘텐츠 선택
  \item \textbf{PDF 인쇄} — 셀별 인쇄 또는 전체 일괄 인쇄 (학생용 시험지 + 정답표)
\end{itemize}

\section*{자주 쓰는 흐름}
\begin{enumerate}
  \item 학생별 탭 → 학생/반 선택 → 매트릭스 표시
  \item 빈 셀 클릭 → ``콘텐츠 배정'' → 검색 모달에서 콘텐츠 선택
  \item 셀 클릭 → 모달 → ``학습 히스토리 보기'' / ``PDF 인쇄''
  \item 일괄 인쇄 — 우측 상단 🖨 → 모든 배정 콘텐츠가 한 PDF 로 (각 학습마다 새 페이지)
\end{enumerate}

\begin{tipbox}{알아두면 좋은 점}
\begin{itemize}[leftmargin=1.2em]
  \item 매트릭스 1행 1열은 행/열 고정 (sticky) — 스크롤해도 헤더가 보입니다
  \item 새 열 추가 시 이름은 자동 생성 (``국어농장 1, 국어농장 2 ...'')
  \item 인쇄 시 학생용 시험지에 \textbf{국어농장 로고 + 본 기관 로고}가 함께 출력됩니다 (기관 설정에서 로고 등록 시)
  \item 콘텐츠 검색 모달에서 \textbf{미리보기} 클릭 시 새 탭이 아닌 in-place 패널로 표시 (세션 격리 방지)
\end{itemize}
\end{tipbox}
